\documentclass[11pt,a4paper]{article}
\usepackage[margin=25mm]{geometry}
\usepackage{amsmath,amssymb,amsthm}
\newtheorem{theorem}{Theorem}
\newcommand{\Sz}{\operatorname{Sz}}
\title{A nonzero coefficient for two Suzuki subfield faces}
\author{Human Compression project: bounded F003 continuation}\date{}
\begin{document}\maketitle
We continue the accepted order-four reduction. Its proper-face,
subfield, cyclic-overlap and compositum arguments remain premises.
The present result removes a degeneracy in a necessary linear equation;
it does not exclude the remaining two-face configuration.

Let $k=\mathbb F_{2^N}$, with $N$ odd, and let
$K=\mathbb F_{2^d}$ and $L=\mathbb F_{2^e}$ be subfields with
$\gcd(d,e)=1$. Write $\theta=2^{(N+1)/2}$, so that
$(z^\theta)^\theta=z^2$ on $k$. We use the matrices $U(x,y)$,
$z(y)=U(0,y)$ and $J$ from the acquired coordinate interface.
The earlier normalization gives
\[
 b=U(1,\epsilon),\quad \epsilon\in\mathbb F_2,\qquad
 u=U(x,0)^{-1}Jz(r)JU(x,0),\qquad
 v=U(w,t)^{-1}Jz(s)JU(w,t),                    \tag{1}
\]
where $x,r\in K$, $w,s\in L$, $r,s\ne0$, and $t\in k$.
Indeed both involutions are away from the unique fixed point of $b$;
the earlier proper-face argument supplies this fact. Central translations
from $K$ and $L$ remove their two original second coordinates, preserving
the condition $t\notin K+L$. Put
$X=x^\theta$, $R=r^\theta$, $W=w^\theta$, and $S=s^\theta$.
The rotation relation is
\[
                 u b^{-1}v=vbu.                         \tag{2}
\]

\begin{theorem}
Suppose (1)--(2) hold and $\langle b,u\rangle$ is a genuine Suzuki
subfield group $\Sz(K)$ with $|K|>2$. Then
\[
                 w^2+\epsilon+X+R/r^2+1\ne0.             \tag{3}
\]
Consequently a necessary equation for $t$, obtained below, has a
nonzero coefficient and determines $t$ uniquely from $x,r,w,s,\epsilon$.
\end{theorem}

\noindent\emph{The matrix identity used in the proof.}
Conjugate (1) by $U(x,0)^{-1}$. The first involution becomes $Jz(r)J$,
the middle rotation becomes $U(1,d)$ with
$d=\epsilon+x+X$, and the last involution has coordinates
$(x+w,t+(x+w)X,s)$. For the following identity use independent symbols
$a,A,d,D$ for $r,R,d,d^\theta$.
Let $V$ be self-adjoint for $J$, so that
\[
 V=\begin{pmatrix}
 v_0&v_1&v_2&v_3\\v_4&v_5&v_6&v_2\\
 v_7&v_8&v_5&v_1\\v_9&v_7&v_4&v_0
 \end{pmatrix}.
\]
Impose $Jz(a)J\,U(1,d)^{-1}V=VU(1,d)Jz(a)J$ and the two trace
equations $\operatorname{Tr}(U(1,d)^2V)=B$ and
$\operatorname{Tr}(U(1,d)V)=Q$. There is a unique solution after setting
$v_1=\sigma$, $v_2=\tau$. The minor using intertwiner positions $(11),(12),(13),(21),(22),(31)$,
the two traces, and all columns except $v_1,v_2$ is $a^2A^2$; the companion file gives every resulting entry.

Set $I=(V^2-I_4)_{13}$, with indices beginning at one, and
$H=c_2(U(1,d)^2V)-B_1$, where $c_2$ is the second characteristic
coefficient. Direct substitution gives the identity
\[
 \Phi=\bigl(1+D(d+D+1+a/A)\bigr)I+DH,                    \tag{4}
\]
where
\begin{align*}
\Phi={}&\sigma\tau+(d+A/a^2)\sigma B
 +(d^2+dA/a^2+1+1/a+A/a^2)\tau B\\
&+(d^3+d^2D+d+d/a+dA/a^2+D^2+D+DA/a^2+A/a^2)B^2\\
&+\sigma Q+(d+a/A+1)\tau Q\\
&+(d^2+da/A+d+D+A/a^2)BQ+DB_1.            \tag{5}
\end{align*}
All denominators in (4)--(5) are powers of $a,A$, which are nonzero.
For an actual involution in the Suzuki group, $I=0$ and $H=0$:
the latter uses the acquired identity $c_2(g)=\operatorname{Tr}(g)^\theta$.
Thus $\Phi=0$ when $B_1=B^\theta$.

\begin{proof}[Proof of the theorem]
For the conjugated involution in (1), matrix multiplication gives
\[
 \begin{split}
 B&=S,\qquad Q=s+S(w^2+w+1),\qquad B_1=s^2,\\
 \tau&=s+S(x+w),\\
 \sigma&=s(X+W)+S\bigl(t+(x+w)X\bigr).
 \end{split}                                                   \tag{6}
\]
Substitute (6), $a=r$, $A=R$, $d=\epsilon+x+X$ and
$D=\epsilon+X+x^2$ into (5). The result is affine in $t$, with coefficient
\[
                    S^2(w^2+\epsilon+X+R/r^2+1).          \tag{7}
\]
Assume it vanishes. Since $w^2\in L$ and the remaining terms lie in $K$,
we have $w^2\in K\cap L=\mathbb F_2$, hence $w\in\mathbb F_2$.
Put $c=\epsilon+1+w$, $p=x+c$, and $P=X+c=p^\theta$.
Equation (7) gives $R/r^2=P$. Both $p,P$ are nonzero.
Applying $\theta$ and using $R^\theta=r^2$ gives
$1= r^2P^2p^2$, hence
\[
                       r=(pP)^{-1},\qquad R=(p^2P)^{-1}.   \tag{8}
\]
If $p=1$, then $x\in\mathbb F_2$ and $r=1$. Both $b$ and $u$ in
(1) then lie in the standard $\Sz(2)$, contrary to the hypothesis.
We may therefore suppose $p\ne1$.

On substituting (8) and $w\in\mathbb F_2$ into the constant term of
(5), the four choices of $\epsilon,w$ give the same equation:
\[
 (p+1)^2s^2+(P+1)^2sS+p^2(P+1)^2S^2=0.                 \tag{9}
\]
Divide by $S^2$ and put $z=s/S\in L^\times$.
This is a nonzero quadratic over $K$, since $p\ne1$.
The degree of $z$ over $K$ is at most two and divides the odd integer
$[k:K]$; therefore $z\in K$. Consequently
$z\in K\cap L=\mathbb F_2$, so $z=1$.
The left side of (9), divided by $S^2$, now equals
$(p+1)^2P^2$, which is nonzero. This contradiction proves (3).
\end{proof}

The accompanying checker performs exact Laurent-polynomial arithmetic
over $\mathbb F_2$. It verifies the rank minor, all matrix and trace
equations, (4), (7), each of the four substitutions giving (9), and the
last nonzero expression. Its monomial inverses record precisely the
denominators used above. No finite ambient group enumeration is a premise.
The remaining task is to prove that the uniquely determined $t$ belongs
to $K+L$, or to find another contradiction. Odd middle rotations with
both face degrees greater than three are also still unresolved. No new
whole-original, conjecture, Lean or measurement credit is claimed.
\end{document}
